% GENERATED by md2tex.py from ../draft_v5_en/body/12_conclusions.md -- do not edit by hand.
% Change the Markdown and run `python md2tex.py`; edits here are overwritten.

\chapter{Conclusions and future work}
\label{ch:12}

\section{Purpose}
\label{sec:12.1}

This chapter closes the work in three parts. It summarises the findings that the thesis supports with evidence, it answers the research questions from Chapter~\ref{ch:01}, and it turns the open issues into concrete lines of work, each with a success criterion. The order follows how much each result is worth, not when it happened, and the negative findings get as much space as the positive ones, because they count just as much as results.

\section{Summary of results}
\label{sec:12.2}

The results are labelled \texttt{R1…R14} and are referred to by that label. This avoids a clash with two sets of labels already used in the document: \texttt{H-01…H-12} are the hazards in the register (Chapter~\ref{ch:04} and Appendix~\ref{app:A}), and \texttt{H1…H3} are the hypotheses stated in \S\ref{sec:1.3} and evaluated in \S\ref{sec:11.3}.

\textbf{R1 — One person can put the full traceability chain into practice.} The chain \texttt{Hazard → Requirement → Rule → Scenario → Metric → Evidence → Verdict} was checked by a tool during the whole project, with no orphans at any review gate, using a checker that cost very little. This is the main claim of the thesis and the one with the strongest evidence.

\textbf{R2 — The literal verdict and the requirement's own criterion can disagree, and the disagreement is useful.} The global verdict of the reference campaign is literally negative. The whole reason is one performance clause in one scenario, while the safety conditions are fully met. Keeping both readings in the record, the literal failure and the argued explanation, is more honest than rewriting the verdict, and it revealed a real defect in a criterion that editing the verdict afterwards would have hidden.

\textbf{R3 — How verdicts are combined matters: indeterminate is not the same as failed.} A per-run verdict that depends on a quantity that was not logged is an instrumentation gap, not a violation. Keeping them apart is what made it possible, once the gaps were closed, to read the new measurements as real results (one positive, one negative) and not as side effects of the aggregator.

\textbf{R4 — The training curve does not tell how the policy behaves.} Seeds with almost identical reward curves produce policies on opposite sides of the line between respecting the constraints and depending on the envelope. In practice, this means the choice has to be made by closed-loop evaluation, with the intervention rate as a main criterion.

\textbf{R5 — The checkpoint at the reward peak can be the worst candidate.} For the reference policy it was: fourteen safety interventions, against zero for the chosen candidate. This follows directly from the previous finding, and it also answers the objection that the policy was chosen after seeing the results, because the criterion was fixed before the campaign.

\textbf{R6 — The cage is latent when the policy respects the constraints, and its value shows where perception gets worse.} With perfect perception the envelope does not intervene inside the domain. With a degraded camera it removes failures that the policy makes on its own. Comparing the two arms is what turns \enquote{the cage adds safety} into a measured claim instead of an assumption about the architecture.

\textbf{R7 — The safety invariant holds: zero road edge contacts inside the domain with the cage active}, against sixty made by the policy without it, over 945 enforcement runs.

\textbf{R8 — End-to-end learning from the camera tracks the lane more precisely than the classical baseline} on a winding layout, which is the opposite of the result on a simple layout. This is the finding that justifies the cost of the learned component in this domain.

\textbf{R9 — Rule composition is a real problem, not a theoretical one.} When the lateral and heading rules fire at the same time, violations of the lateral margin still happen. Better training halves the problem but does not change its nature, which makes it a design problem of the envelope. It is the only requirement that the work closes as not satisfied, and it is reported that way.

\textbf{R10 — Fail-safe caution costs availability, and the cost can be measured.} The controlled stops that prevent excursions also cut short runs that would have ended well. The design of the trigger is an explicit trade-off between safety and availability, not a free choice.

\textbf{R11 — The gap between simulators is a result in itself.} Moving the system to a more realistic environment meant recalibrating the dynamic authority, the heading thresholds and even the structure of the reward. Checkpoints do not transfer. That negative result set expectations for the move to hardware better than any optimistic estimate, and the hardware confirmed it. The policy validated in simulation keeps the correct sign of its response to the lane, but its amplitude is ten times too small, against a constant bias 2.1 times larger than the whole response. The cause was the handedness of the training circuit, memorised as a steering bias. Driving needed a retraining aimed at transfer, not a small adjustment.

\textbf{R12 — The cage does not only filter the policy: it shapes it.} Trained with the envelope in the actuation chain, the reference policy produces an almost all-or-nothing command that saturates the rate limiter in 77.5\,\% of the steps. With the cage active it drives better than any other configuration in this work. Without it, it leaves the lane about every three laps. What has been built, and what was taken to hardware, is the pair, not the policy. The dependency is measured, and its origin is inferred. The transfer of the pair has been measured once: with the checkpoint chosen for independence from the cage instead of by reward, the limiter acts in 3.4\,\% of the cycles on the real vehicle, against 3.0\,\% in simulation. The risk was stated before measuring, and the first measurement does not confirm it. With one run and no scored scenario this does not close the question. The physical campaign is what will decide.

\textbf{R13 — An assumption that the simulation inherits is an assumption the simulation cannot falsify.} The camera's effective field of view was a default configuration value that the simulator copied, and every lateral quantity the cage uses depended on it. The first measurement on hardware showed it was wrong (77.89° against the assumed 90°). On the unrectified image, the lane boundary threshold, set at 160 mm, could fire with the vehicle at a real 207–241 mm. Rectifying the image to the standard camera model removed the defect: when the measurement was repeated on the deployed path, the scale is 1.058/0.991 with no intercept, and the threshold fires at a real 151–158 mm. The correction of that second number is part of the finding and does not weaken it: neither the inherited error nor the effect of its fix could be seen in simulation, because the validation loop was closed on itself. It is the strongest argument in the work for making the physical step required and not optional.

\textbf{R14 — Some defects only appear with a real vehicle in front of you.} These are rules that are correct according to their specification but whose behaviour \emph{in operation} is not defined. One example is an emergency stop whose asymmetric exit is correct for a simulated episode, which ends, but leaves a real vehicle that has to keep running stopped forever. Others are thresholds that cannot fire on commanded motion in the deployed configuration, and sensor failures that go straight into the cage's only speed input. None of these is a failure of the policy or of the cage. They are gaps between an artefact validated on episodes and a vehicle that runs continuously.

\section{Answers to the research questions}
\label{sec:12.3}

\textbf{Main question.} \emph{Can the V-Model be adapted with a limited, traceable set of changes so that it can include components trained with reinforcement learning, without losing the two-way link between specification and V\&V?} Yes, based on the evidence of one case. Five adaptations were enough to cover every failure mode found. None of them required giving up the structure of the V. The two-way link was not only kept but made stronger, going from a recommendation to a constraint checked by a tool. The limit of this answer is that it comes from one case: it shows that the adaptations are enough in practice, not that they are complete.

\textbf{Second question.} \emph{Does the framework produce consistent and traceable evidence, including an honest description of the sim-to-real gap?} Yes for traceable evidence, and partly for the gap. All fourteen requirements have a supported verdict, including the unfavourable ones. The description of the gap is structured: its first step is measured and its second step has gone through bring-up. The high-fidelity bridge gave a useful negative diagnosis. The physical platform, in the bring-up phase, produced a falsified assumption, a negative result about the validated policy, a positive preliminary observation about its replacement, the first measurement of a transfer risk stated in advance, and, measured with true-position methods, the finding that the lane estimator's accuracy depends on the place on the circuit and not on what the vehicle does. No simulation campaign could have shown that failure mode, because in simulation the estimator is never in a real place. What has not been produced is the physical campaign, and so there is no per-scenario verdict on hardware. The measurement exists as bring-up and not as a verdict, the physical column of the final table is still empty for that reason, and this is stated openly. The part of A5 that is still open is therefore smaller than before and is described precisely instead of as a simple absence, but it is still open.

\section{Future work}
\label{sec:12.4}

\textbf{T1 — Designing and verifying how rules are combined.} This is the open problem that no policy improvement fixed, and the only unmet requirement. The next step is to move from combining rules in a fixed order to analysing the joint envelope. That means finding the region of the state space where applying the rules one after another gives actions outside the envelope that each rule guarantees on its own, and then either reordering and re-tuning the rules, or recording that region as excluded from the domain. The breakdown by rule combination already shows where to look (the lateral limit and heading rules firing together), and the existing scenario grid is the tool to measure it.

\textbf{T2 — Scoring scenarios on hardware.} The physical deployment is done and the vehicle drives. What is missing is the half that produces a verdict. Four conditions stand in the way, and all four are known. One of them was measured after the vehicle drove for the first time: widen the lane estimator, whose accuracy depends on the place on the circuit and not on the action, before restricting the policy. The other three are: a reset path for the emergency stop, so a scored run does not end at the first perception pulse (solved outside the cage, precisely to avoid changing the artefact under verification); per-run provenance logging, which is already implemented; and a decision on which perception contract to use, because the heading setting that lets the vehicle drive is not the one the simulation campaigns were scored with. This work predicted which component would transfer best: the cage, because it is specified on an abstract state. That prediction held. No safety rule fired during the run, and the only rule that acted did so at the rate simulation had predicted.

\textbf{T3 — Ablation of the rate limiter.} The measurement behind R12 is done, but its mechanism is still inferred. Retraining with the limiter outside the actuation loop and comparing the distribution of the raw command would answer the question, and it is cheap to do in simulation. The original advice was to do this before moving to hardware. Hardware came first, but that does not make the ablation unnecessary: the fact that the pair transfers at the expected rate shows that the dependency survives a change of plant, not where it comes from.

\textbf{T4 — Robustness of the cage estimator over time.} There is no reliable single-frame fix for a reading that is wrong but consistent with itself. The directions with a solid basis are: consistency over several frames (a stable false reading in the same section of the circuit would show up as a mismatch with the odometry), calibrating the estimator's confidence against the simulation's ground truth, and cheap sensor redundancy for the cage state.

\textbf{T5 — Availability: a more robust emergency trigger.} Tune hysteresis and persistence in the validity and perception health triggers so that false stops do not happen, without making the reaction to real degradation slower. The acceptance criterion can already be measured with the existing library.

\textbf{T6 — Checking execution, not only references.} This is the improvement the framework found in itself. Extending the checker to verify that a scenario runs as specified (that its starting conditions are injected and that the quantities its criterion uses are logged) would close the gap that let one scenario run for months without doing what it claimed.

\textbf{T7 — Replication and scale.} Apply the framework to a second case with a learned component and on a photorealistic simulator, to separate what belongs to the framework from what belongs to this case and to this tool.

\section{Conclusion}
\label{sec:12.5}

The thesis argues that a classical functional safety life cycle can include a component trained with reinforcement learning without losing what makes it valuable, which is the verifiable link between what is specified and what is verified. It supports this with a complete case whose chain of evidence can be audited from start to finish.

The most solid result is also the simplest. Inside the stated operational domain, with the envelope active, the system did not leave the road once in nine hundred and forty-five runs, while without it, it did so sixty times. The most useful result, however, is uncomfortable: the envelope that makes that number possible also shaped the policy inside it, to the point where the policy can no longer work without it. And the result the work leaves as an open problem is the only requirement it does not meet: how simultaneous rules are combined, measured on two different policies and still present in both.

There is a fourth result that does not come from simulation, and that the work would not have reached without simulation either. The physical step did not confirm the simulation results. It did not score a single scenario, and that is stated openly. What it did was falsify assumptions that simulation could not reach: a camera parameter that the simulator had inherited and therefore copied, the handedness of a training circuit learned as a steering bias, and an estimator whose accuracy depends on where the vehicle is and not on what it does. No amount of extra simulation would have revealed any of the three, because in all three cases the measuring tool shared the error with the thing being measured. This is the strongest argument this thesis can give for making the measurement of the gap a required part of the cycle and not an optional appendix, and it is also why adaptation A5 is written as a requirement and not as a recommendation.

These four results sit in the same document, with the same traceability, and none of them was rewritten to look better. That is the best evidence available that the framework does what it says it does.
